\selectlanguage{spanish}

\renewcommand{\articuloidiomaxml}{es}
\renewcommand{\articulotitulo}{Reseña bibliográfica}
\renewcommand{\articuloautores}{Valentina Polack}
\renewcommand{\articulotipo}{RESEÑA}
\renewcommand{\articulorevista}{Revista Valen de ideas}
\renewcommand{\articuloissne}{0262-2130}
\renewcommand{\articulovolumen}{01}
\renewcommand{\articulonumero}{01}
\renewcommand{\articulofecha}{ 2026}

\setcounter{page}{1}
\thispagestyle{firstpage}

% --- TÍTULO A ANCHO COMPLETO ---
\noindent\begin{minipage}[t]{\dimexpr\textwidth+\marginparsep+\marginparwidth\relax}%
{\sffamily\bfseries\Large\articulotitulo\par}%
\end{minipage}\par
\vspace{3mm}%
\renewcommand{\articulotitulotrans}{Book review}
\noindent\begin{minipage}[t]{\dimexpr\textwidth+\marginparsep+\marginparwidth\relax}%
{\sffamily\itshape\normalsize\articulotitulotrans\par}%
\end{minipage}\par
\vspace{4mm}%

% --- BLOQUE LATERAL DE METADATOS ---
\begin{textblock*}{68mm}(\gbbloquex,92mm)
{\setlength{\leftskip}{0pt}\setlength{\parindent}{0pt}%
\noindent\begin{minipage}[t]{68mm}
\sffamily\small\setlength{\parskip}{1pt}
{\bfseries\color{azulrevista}Valentina Polack}\par
{\scriptsize\texttt{\href{https://orcid.org/0009-0007-5867-5325}{https://orcid.org/0009-0007-5867-5325}}}\par
\vspace{6pt}
{\bfseries\color{azulrevista!70}Derechos de autor\'ia:}\par
\textcopyright\ 2026 Polack, V. Publicado por Revista Valen de ideas. Se permite el uso, distribución y reproducción sin restricciones, incluso con fines comerciales, siempre que se cite la obra original. (\href{https://creativecommons.org/licenses/by/4.0/}{https://creativecommons.org/licenses/by/4.0/})\par\vspace{6pt}
{\bfseries\color{azulrevista!70}Publicación:} 2026, vol.~01, n\'um.~01\par\vspace{6pt}
{\bfseries\color{azulrevista!70}URL:} {\scriptsize\texttt{\href{https://www.wattpad.com/}{https://www.wattpad.com/}}}\par\vspace{6pt}
{\bfseries\color{azulrevista!70}Licencia:} Creative Commons Attribution (CC BY 4.0)\par
{\scriptsize\texttt{\href{https://creativecommons.org/licenses/by/4.0/}{https://creativecommons.org/licenses/by/4.0/}}}\par
\vspace{6pt}
\vspace{6pt}
{\bfseries\color{azulrevista!70}Compartir:}\par\vspace{3pt}
{\large
\href{https://bsky.app/intent/compose?text=https%3A%2F%2Fwww.wattpad.com%2F}{\includesvg[height=0.9em]{/home/valentina/.gbpublisher/fonts/bluesky}}~
\href{mailto:?subject=Rese%C3%B1a%20bibliogr%C3%A1fica&body=https%3A%2F%2Fwww.wattpad.com%2F}{\includesvg[height=0.9em]{/home/valentina/.gbpublisher/fonts/envelope}}~
\href{https://www.facebook.com/sharer/sharer.php?u=https%3A%2F%2Fwww.wattpad.com%2F}{\includesvg[height=0.9em]{/home/valentina/.gbpublisher/fonts/facebook}}~
\href{https://www.linkedin.com/sharing/share-offsite/?url=https%3A%2F%2Fwww.wattpad.com%2F}{\includesvg[height=0.9em]{/home/valentina/.gbpublisher/fonts/linkedin}}~
\href{https://www.reddit.com/submit?url=https%3A%2F%2Fwww.wattpad.com%2F&title=Rese%C3%B1a%20bibliogr%C3%A1fica}{\includesvg[height=0.9em]{/home/valentina/.gbpublisher/fonts/reddit}}~
\href{https://www.researchgate.net/search?q=Rese%C3%B1a%20bibliogr%C3%A1fica}{\includesvg[height=0.9em]{/home/valentina/.gbpublisher/fonts/researchgate}}~
\href{https://t.me/share/url?url=https%3A%2F%2Fwww.wattpad.com%2F&text=Rese%C3%B1a%20bibliogr%C3%A1fica}{\includesvg[height=0.9em]{/home/valentina/.gbpublisher/fonts/telegram}}~
\href{https://wa.me/?text=https%3A%2F%2Fwww.wattpad.com%2F}{\includesvg[height=0.9em]{/home/valentina/.gbpublisher/fonts/whatsapp}}~
\href{https://twitter.com/intent/tweet?url=https%3A%2F%2Fwww.wattpad.com%2F&text=Rese%C3%B1a%20bibliogr%C3%A1fica}{\includesvg[height=0.9em]{/home/valentina/.gbpublisher/fonts/x-twitter}}
}\par
\end{minipage}%
}% FIN GRUPO leftskip
\end{textblock*}


\bigskip
\noindent\rule{\linewidth}{0.4pt}
\bigskip


\begin{quote}
Souza, Jessé. \textit{A elite do atraso: da escravidão à Lava Jato}. São Paulo: Leya, 2017, 242 pp. ISBN 978-85-441-0408-1.

\end{quote}


\section{Presentación y contexto}
Jessé Souza es uno de los sociólogos brasileños más provocadores de su generación. Formado en Derecho y Sociología por la Universidad de Brasilia, doctorado en Heidelberg y con posdoctorados en la New School for Social Research de Nueva York, ha dedicado buena parte de su trayectoria a desmontar lo que considera los mitos fundacionales de la interpretación de Brasil como nación. \textit{A elite do atraso}, publicado originalmente en 2017 en el contexto de la crisis político-institucional que desembocó en el \textit{impeachment} de Dilma Rousseff, es quizás su obra de mayor repercusión pública: un libro que trascendió el circuito académico para instalarse en el debate político brasileño e incluso inspirar el desfile de la escuela de samba Paraíso do Tuiuti en el carnaval de Río de Janeiro de 2018.

El argumento central del libro puede resumirse en una operación de desplazamiento: donde la tradición dominante del pensamiento social brasileño ubicó al patrimonialismo heredado de Portugal como la clave explicativa de los males nacionales —la corrupción, el atraso institucional, la ineficiencia del Estado—, Souza propone que la verdadera matriz de la desigualdad brasileña es la esclavitud y su legado persistente. No se trata, advierte el autor, de una tesis meramente historiográfica: lo que está en juego es la estructura de legitimación que permite a las élites brasileñas contemporáneas reproducir privilegios mientras trasladan la responsabilidad del “atraso” al Estado y, por extensión, a las clases populares que dependen de él.

Para sostener esta tesis, Souza emprende una revisión crítica de tres figuras canónicas del pensamiento social brasileño: Gilberto Freyre, Sérgio Buarque de Holanda y Raimundo Faoro. El argumento no es que estos autores carecieran de mérito intelectual, sino que sus interpretaciones —el “hombre cordial”, el patrimonialismo ibérico, la confusión entre lo público y lo privado— terminaron configurando un sentido común que opera ideológicamente a favor de los intereses de clase de la élite. Al atribuir los problemas de Brasil a una herencia cultural portuguesa genérica, estas narrativas invisibilizan el papel de la esclavitud como institución estructurante de las relaciones sociales y, sobre todo, desplazan la atención de la explotación económica real hacia supuestos defectos culturales.

El libro se organiza en torno a una doble operación. La primera es demoledora: Souza desmonta con considerable vigor polémico la noción de que el patrimonialismo constituya una categoría explicativa adecuada para comprender la sociedad brasileña. Según el autor, la tesis patrimonialista produce un efecto de ocultamiento al sugerir que la corrupción es un problema fundamentalmente estatal, lo que permite eximir de responsabilidad a las formas de acumulación privada que históricamente se beneficiaron de la explotación de mano de obra esclavizada y, posteriormente, de la precarización del trabajo libre. La segunda operación es propositiva: Souza ofrece una lectura alternativa de la formación social brasileña que pone en el centro la continuidad entre la esclavitud colonial y las formas contemporáneas de desprecio social hacia los pobres.

Uno de los aportes más sugerentes del libro es la noción de “ralé” (\textit{canalla} o \textit{chusma}) como categoría sociológica. Souza argumenta que existe una clase social estructuralmente excluida —no simplemente pobre en términos de ingreso— cuya exclusión se reproduce a través de mecanismos que no son puramente económicos sino también simbólicos y afectivos. El desprecio hacia esta clase, naturalizado en la vida cotidiana brasileña, constituye para Souza la herencia más duradera y menos reconocida de la esclavitud.

Desde una perspectiva latinoamericana más amplia, el libro plantea preguntas que exceden el caso brasileño. La relación entre élites económicas, producción intelectual y legitimación del orden social es un problema compartido por todas las sociedades de la región. ¿En qué medida las ciencias sociales latinoamericanas han contribuido, inadvertidamente, a naturalizar las desigualdades que pretenden explicar? ¿Cuánto de lo que se presenta como análisis desinteresado de la cultura nacional opera, en realidad, como justificación del privilegio?

El libro no está exento de debilidades. La crítica a Freyre, Buarque de Holanda y Faoro, aunque estimulante, resulta por momentos esquemática: al subordinar la lectura de estos autores a la demostración de su tesis, Souza tiende a simplificar obras de considerable complejidad. La categoría de “ralé”, por su parte, si bien tiene poder descriptivo, plantea problemas analíticos que el libro no resuelve del todo: ¿se trata de una clase en sentido estricto o de una condición social? ¿Cuáles son sus fronteras empíricas? Finalmente, el tono polémico que recorre toda la obra —y que sin duda contribuyó a su éxito editorial— puede resultar excesivo para un público académico habituado a la argumentación matizada.


\section{Conclusión}
Estas objeciones, sin embargo, no disminuyen el valor del libro como intervención intelectual. \textit{A elite do atraso} cumple con la función que toda buena obra de ciencias sociales debería cumplir: obligar a repensar categorías que se daban por sentadas. Para los lectores de habla hispana, el texto ofrece además una oportunidad de reflexionar sobre las propias tradiciones interpretativas nacionales y sobre el modo en que las élites de cada país han construido relatos funcionales a la reproducción de sus privilegios. En un momento en que las derechas radicales avanzan en toda la región apelando a discursos que combinan la meritocracia individual con el desprecio hacia los sectores populares, la lectura de Souza resulta no solo pertinente sino urgente.


